
Nous détaillerons l'architecture du programme C de notre version de initiale de \textit{Royal Dédale}
\begin{figure}[H]
	\centering
	\includegraphics[width=0.5\textwidth]{architecture_v0.png}
	\caption{Architecture de la version initiale de notre projet}
	\label{fig:architecture_v0}
\end{figure}

La figure \ref{fig:architecture_v0} nous montre les modules de nos programme et leurs interdépendances, chaque flèche indique que le fichier au bout de la flèche est dépendant de celui à l'origine mais aussi de tout les modules dont ce dernier en dépend. Nous pouvons regrouper les modules en trois groupes:
\begin{itemize}
  \item La \textbf{fondation}. Où on implémente les "types" (structures et énumérations) fondamentales pour notre programme: $localisation$, $direction (dir)$, $ressource$ et $stash (poche)$; et leurs opération associées. Ces conceptes servent de base à la quel s'appuie les modules suivents.
  \item L'\textbf{implémentation des tuiles}. Avec les modules précédents nous créeons les tuiles et leurs opérations. C'est de ça qui est composé le manoire implémenter par la suite et une mécanique fondamentale de la partie. 
  \item  L'\textbf{implémentation du jeu}. C'est la que nous créons le manoir, les fonctions pour le manipuler, des fonction pour l'afficher en forme d'interface textuelle et combinons tout ça dans une boucle qui nous permettra de jouer le jeu. Dans cette étape nous combinons tout notre code précédant pour créé finalment le jeu.
\end{itemize}

\vspace{1em}
Cet architecture donne un jeu fonctionnel mais forcément de la manière la plus optimal ou complet en terme de jouabilité, nous verrons les limites et les amélioratoins proposé à ce programme. Mais avant ça, nous commençons par detailler la fondation de cette architecture.
